\documentclass[11pt]{article}
\usepackage[T1]{fontenc}
\usepackage[ngerman]{babel}
\usepackage{amsmath,amssymb,amsthm}
\usepackage[a4paper,margin=2.6cm]{geometry}
\usepackage{hyperref}
\usepackage{lmodern}
\usepackage{microtype}

\title{Eine spiegelsymmetrische Polyederschalen-Konstruktion\\
	eines dodekaederartigen Zwölfflächners}
\author{}
\date{}

\newtheorem{satz}{Satz}
\newtheorem{proposition}[satz]{Proposition}
\newtheorem{definition}[satz]{Definition}
\newtheorem{bemerkung}[satz]{Bemerkung}

\begin{document}
	\maketitle
	
	\begin{abstract}
		Wir untersuchen eine in einer schematischen Tafel dargestellte Polyederkonstruktion, die aus zwei kongruenten, spiegelsymmetrisch gepaarten Halbschalen besteht. Jede Halbschale wird durch ein planares Netz mit ausgezeichneten Faltkanten beschrieben und besitzt eine hexagonale Kontaktfläche, entlang derer zwei Hälften zu einem geschlossenen Körper zusammengesetzt werden. Die resultierende Figur wird in der Vorlage als ``special case of a halved dodecahedron'' bezeichnet. Wir zeigen, dass diese Bezeichnung im strengen Sinn irreführend ist: Der entstehende Körper ist im Allgemeinen kein halbiertes reguläres Dodekaeder und auch nicht notwendig durch eine ebene Halbierung eines regulären Dodekaeders entstanden. Stattdessen handelt es sich um eine eigene nicht-reguläre polyedrische Konstruktion mit zwölf äußeren Flächen und spiegelsymmetrischer Montage.
	\end{abstract}
	
	\section{Einleitung}
	
	Polyedrische Netze und deren Faltungen bilden ein klassisches Thema der diskreten Geometrie. Von besonderem Interesse sind Konstruktionen, bei denen aus ebenen polygonalen Teilfiguren durch Faltung räumliche Schalen entstehen, die sich anschließend entlang gemeinsamer Randflächen zu einem geschlossenen Polyeder zusammensetzen lassen.
	
	Die vorliegende Abbildung beschreibt genau eine solche Situation: Ein annotiertes Netz definiert eine räumliche Halbschale; zwei identische Exemplare, von denen eines gespiegelt ist, werden entlang einer hexagonalen Grenzfläche zusammengesetzt. Der entstehende Körper besitzt nach der Darstellung zwölf äußere Flächen und soll dodekaedralen Charakter haben.
	
	\section{Grundbegriffe}
	
	\begin{definition}
		Ein \emph{reguläres Dodekaeder} ist ein konvexes Polyeder mit genau zwölf Flächen, wobei jede Fläche ein reguläres Pentagon ist und an jedem Eckpunkt genau drei Flächen zusammentreffen.
	\end{definition}
	
	\begin{definition}
		Eine \emph{Polyederschale} sei ein aus polygonalen Flächen zusammengesetzter, zusammenhängender polyedrischer Teilkörper mit nichtleerem Rand.
	\end{definition}
	
	\begin{definition}
		Eine \emph{Kontaktfläche} sei eine zusammenhängende polygonale Randfläche, entlang derer zwei Polyederschalen kantenkongruent identifiziert werden können.
	\end{definition}
	
	\begin{definition}
		Ein \emph{dodekaederartiger Zwölfflächner} sei hier informell ein Polyeder mit zwölf äußeren Flächen, ohne dass Regularität, Flächengleichheit oder Pentagonalität vorausgesetzt werden.
	\end{definition}
	
	\section{Warum die Konstruktion kein halbiertes reguläres Dodekaeder ist}
	
	Die vorgelegte Grafik enthält Winkelangaben wie $30^\circ$, $90^\circ$ und $120^\circ$, ferner eine ausdrücklich markierte hexagonale Kontaktfläche sowie Hinweise auf nicht-reguläre Flächen. Dies steht im deutlichen Kontrast zum regulären Dodekaeder.
	
	\begin{satz}
		Die in der Tafel dargestellte Konstruktion beschreibt im Allgemeinen keine Halbierung eines regulären Dodekaeders.
	\end{satz}
	
	\begin{proof}
		Die Konstruktion verwendet nicht-reguläre Flächen, eine hexagonale Kontaktfläche und lokale Winkelkonfigurationen, die nicht mit einer Zerlegung des regulären Dodekaeders in zwei kongruente, entlang einer einfachen Halbierungsfläche getrennte Hälften identisch sind. Insbesondere folgt aus der expliziten Kennzeichnung ``non-regular faces'', dass keine reguläre dodekaedrische Flächenstruktur vorliegt. Damit kann der Endkörper allenfalls dodekaederartig, nicht aber regulär-dodekaedrisch sein.
	\end{proof}
	
	\section{Was die Konstruktion tatsächlich liefert}
	
	Man konstruiert eine einzelne nicht-reguläre Halbschale mit hexagonaler Grenzfläche. Zwei kongruente Exemplare werden spiegelbildlich zueinander orientiert und entlang der gemeinsamen Fläche verklebt.
	
	\begin{proposition}
		Unter der Annahme, dass die Faltung des Netzes widerspruchsfrei ist und die Kontaktflächen isometrisch übereinstimmen, ergibt die Montage zweier gespiegelter Hälften einen geschlossenen polyedrischen Körper mit einer Spiegelebene entlang der Naht.
	\end{proposition}
	
	\begin{bemerkung}
		Dies ist eine strukturelle, keine metrische Aussage. Sie setzt weder Konvexität noch Regularität voraus.
	\end{bemerkung}
	
	\section{Notwendige Konsistenzbedingungen}
	
	Damit die Konstruktion mathematisch vollständig legitimiert ist, müssen mehrere Bedingungen überprüft werden:
	\begin{enumerate}
		\item planare Konsistenz des Netzes,
		\item globale Faltbarkeit ohne Selbstüberschneidung,
		\item Randkongruenz der Kontaktflächen,
		\item Eck- und Kantenverträglichkeit an der Naht,
		\item topologischer Typ des Endkörpers.
	\end{enumerate}
	
	\section{Terminologische Korrektur}
	
	Der Ausdruck
	\[
	\text{``a special case of a halved dodecahedron''}
	\]
	ist mathematisch irreführend. Präziser wären Formulierungen wie:
	\begin{itemize}
		\item spiegelsymmetrisch gepaarte Polyederschalen-Konstruktion,
		\item nicht-regulärer dodekaederartiger Zwölfflächner,
		\item durch hexagonale Kontaktfläche montierter Zwölfflächner.
	\end{itemize}
	
	\section{Schluss}
	
	Die Tafel zeigt eine geometrisch interessante und didaktisch gut visualisierte Polyederidee. Ihre Stärke liegt in der klaren Darstellung des Prinzips
	\[
	\text{Netz} \longrightarrow \text{Halbschale} \longrightarrow \text{Spiegelpaarung} \longrightarrow \text{Gesamtkörper}.
	\]
	Ihre Schwäche liegt in der Terminologie. Korrekt beschrieben handelt es sich um eine nicht-reguläre, spiegelsymmetrische Polyederkonstruktion mit zwölf Außenflächen, nicht um ein halbiertes Dodekaeder im klassischen Sinn.
	
\end{document}